% Part: incompleteness
% Chapter: theories-computability
% Section: q-is-ce

\documentclass[../../../include/open-logic-section]{subfiles}

\begin{document}

\olfileid{inc}{tcp}{qce}

\olsection{$\Th{Q}$ iku \printtoken{S}{c.e.}-jangkep}


\begin{thm}
$\Th{Q}$ iku !!{c.e.} nanging ora bisa diputusake. Malah,
  iki himpunan !!{c.e.} kang jangkep.
\end{thm}

\begin{proof}
Ora angel kanggo weruh yen $\Th{Q}$ iku !!{c.e.}, amarga iki himpunan
(kode kanggo) ukara $y$ kang nduweni bukti $x$ tumrap $y$ ing
$\Th{Q}$:
\[
Q = \Setabs{y}{\lexists[x][\Prf[\Th{Q}](x,y)]}.
\]
Nanging kita ngerti yen $\Prf[\Th{Q}](x,y)$ komputabel (malah rekursif
primitif), lan saben himpunan kang bisa ditulis kaya ing ndhuwur iku c.e.

Nyebut iki himpunan c.e.\ kang jangkep padha karo nyebut $K
\leq_m Q$, ing ngendi $K = \Setabs{x}{!A_x(x) \downarrow}$. Mula ayo
dituduhake yen $K$ bisa direduksi menyang $\Th{Q}$. Amarga predikat
Kleene $T(e,x,s)$ rekursif primitif, predikat iki bisa diwakili
ing $\Th{Q}$, upamane dening $!A_T$. Mula kanggo saben $x$, kita nduweni
\begin{align*}
x \in K & \lif \lexists[s][T(x,x,s)] \\
& \lif \lexists[s][(\Th{Q} \Proves !A_T(\num x,\num x, \num s))]
  \\
& \lif \Th{Q} \Proves \lexists[s][!A_T(\num x, \num x, s)].
\end{align*}
Kosok baline, yen $\Th{Q} \Proves \lexists[s][!A_T(\num x, \num x, s)]$,
mula, amarga kabeh aksioma Q bener ing model baku wilangan natural lan aturan bukti njaga bebener, kanggo sawijining wilangan natural $n$, formula $!A_T(\num x,
\num x, \num n)$ mesthi bener. Saiki, yen $T(x,x,n)$ salah,
$\Th{Q}$ bakal mbuktekake $\lnot !A_T(\num x, \num x, \num n)$, amarga
$!A_T$ makili $T$. Nanging banjur $\Th{Q}$ mbuktekake
formula kang salah, lan iki kontradiksi. Mula $T(x,x,n)$ mesthi bener,
kang ateges $!A_x(x) \downarrow$.

Cekake, kanggo saben $x$, $x$ kalebu ing $K$ yen lan mung yen
$\Th{Q}$ mbuktekake $\lexists[s][!A_T(\num x,\num x,s)]$. Dadi fungsi $f$
kang nggawa $x$ menyang (kode kanggo) ukara $\lexists[s][!A_T(\num x,
  \num x, s)]$ iku reduksi saka $K$ menyang $\Th{Q}$. Cathetan koreksi sumber OLPL-791: simpulan lan ukara kang dikodeake saiki nganggo formula representasi A-T, dudu predikat Kleene T ing tataran meta.
\end{proof}

\end{document}
